\section{Results}

\subsection{Overview and evidence tiers}
\label{sec:results-overview}

Results are organized in three parts aligned with the evidence tiers in Table~\ref{tab:objective-evidence-map}.
\textbf{Part~A} (\S\ref{sec:results-assembly-preprocessing}) presents primary evidence on implementation auditability and hydrography preprocessing (Objectives~1--2).
\textbf{Part~B} (\S\ref{sec:results-cal-val-sensitivity}) presents secondary evidence on structural scaling and assembly cost (Objective~3) together with downstream workflow readiness on controlled gage basins (Objective~4).
\textbf{Part~C} (\S\ref{sec:results-runtime-benchmark}) presents supporting evidence on simulation runtime (Objective~5).
The small, medium, and large benchmark domains illustrate discretization scaling; controlled gage basins test workflow execution; the runtime catalog reports simulation cost on one host/build.

\subsection{Assembly and preprocessing evidence}
\label{sec:results-assembly-preprocessing}

\subsubsection{Objective 1: Implementation and workflow architecture}
\label{sec:results-obj1}

The automated workflow produces a standard SWAT\textsuperscript{+} project tree---vector hydrography, gridded and station climate inputs, SQLite project tables, and optional QA sidecars---through the layered pipeline in Figure~\ref{fig:workflow}.
Asynchronous worker processes (Figure~\ref{fig:architecture-layers}) separate long-running model assembly from interactive use so outputs remain in a predictable per-basin directory layout that can be inspected and audited.
Table~\ref{tab:package-manifest-abridged} summarizes representative output file classes; complete path patterns are listed in Supplementary Table~S1.

Structural footprint metrics for the S/M/L benchmark size classes appear in Table~\ref{tab:benchmark-sm-l-products} (\S\ref{sec:results-obj3}).

\subsubsection{Objective 2: Hydrography preprocessing outcomes}
\label{sec:results-obj2}

National NHDPlus HR vectors are transformed into SWAT\textsuperscript{+}-compatible channel, subbasin, and lake layers through the staged pipeline in Figure~\ref{fig:nhd-workflow}.
Table~\ref{tab:nhd-rules-abridged} lists core rules applied before model-scale clipping; the complete rule catalog is in Supplementary Table~S1.
Table~\ref{tab:nhd-preprocessing-qa} compares original USGS NHDPlus HR (HU4 geodatabase, domain-clipped) with exported SWAT\textsuperscript{+} shapefiles from the generated projects.
For the small (S) benchmark domain, 56 original catchments and 56 flowlines include two divergence-code-2 reaches and two isolated reaches (two isolated catchments merged); four orphan catchments are dissolved before the workflow output retains 53 SWAT\textsuperscript{+} watershed polygons and 45 routed channels.
The medium (M) domain shows eleven divergence-code-2 drops and eleven orphan catchments among 1{,}382 original catchments (no isolated reaches in domain), ending at 1{,}371 final watershed polygons and channels.
The large (L) domain includes 88 isolated reaches, 282 orphan catchments, and 34 no-VAA catchments among 9{,}739 original catchments; final routed counts are 8{,}181 channels and 162 subbasins per the structural inventory.
Input counts are recomputed from USGS HU4 geodatabase archives; isolated and orphan catchment columns apply SWATGenX stream-drop rules in memory on the domain-clipped original graph before comparing to exported vectors.
These counts document structural rule application, not demonstrated hydrologic benefit from each preprocessing choice.

\subsubsection{Objective 2 (continued): Drainage-area fidelity audit}
\label{sec:results-drainage-audit}

Structural counts confirm that preprocessing rules were applied; they do not confirm that the \emph{executable} model preserves cumulative drainage-area structure.
To close that gap, drainage area from each generated SWAT\textsuperscript{+} project's \texttt{chandeg.con} was compared with original NHDPlus HR \texttt{TotDASqKm} at gage-assigned channels (Section~\ref{sec:methods-drainage-audit}) across all eight catalog Model IDs (Table~\ref{tab:drainage-area-audit}).
Of 98 audited gages, 97 matched a SWAT\textsuperscript{+} channel and 82 (85\%) fell within the loose $0.5$--$2.0\times$ agreement band; the single unmatched gage lacked a \texttt{chandeg.con} row (a reservoir/auxiliary channel).
Per-model median SWAT\textsuperscript{+}/NHD ratios span $0.57$--$1.16$, with the lowest medians driven by a small number of gage--reach assignment outliers on the sparsely gaged scaling domains (Upper Gila X60, Little Kanawaha X20) rather than by basin-wide conversion error.

\begin{table}[htbp]
  \centering
  \caption{Drainage-area fidelity audit: SWAT\textsuperscript{+} \texttt{chandeg.con} area versus original NHDPlus HR \texttt{TotDASqKm} at gage-assigned channels, by catalog Model ID (Objective~2). ``Outliers'' are gage--channel assignment outliers reported separately from the agreement band.}
  \label{tab:drainage-area-audit}
  \footnotesize
  \setlength{\tabcolsep}{4pt}
  \renewcommand{\arraystretch}{1.05}
  \resizebox{\textwidth}{!}{%
    % Generated by publication/analysis/scripts/emit_tab_drainage_area_audit_tex.py
% Source: /data/SWATGenXApp/codes/web_application/frontend/src/data/drainageAreaAuditCatalog.json (lastUpdated 2026-06-02)
\begin{tabular}{@{}lllrrrrr@{}}
\toprule
Tier & Catalog ID & State & Stations & Matched & Median SWAT/NHD & Within 0.5--2.0$\times$ & Outliers \\
\midrule
S & \texttt{03080102} & FL & 3 & 3 & 0.85 & 3 & 0 \\
M & \texttt{09471300} & AZ & 3 & 3 & 1.16 & 3 & 0 \\
L & \texttt{03100101} & FL & 76 & 76 & 1.11 & 66 & 6 \\
X20 & \texttt{03152000} & WV & 6 & 6 & 0.84 & 4 & 2 \\
X40 & \texttt{07174000} & KS & 3 & 2 & 1.00 & 2 & 0 \\
X60 & \texttt{15060105} & AZ & 4 & 4 & 0.57 & 2 & 2 \\
Cal & \texttt{02297600} & FL & 1 & 1 & 1.08 & 1 & 0 \\
Cal & \texttt{05536265} & IL & 2 & 2 & 0.88 & 1 & 1 \\
\midrule
\textbf{All} & --- & --- & \textbf{98} & \textbf{97} & --- & \textbf{82} & \textbf{11} \\
\bottomrule
\end{tabular}
%
  }
\end{table}

\paragraph{USGS station--channel assignment (Peace River).}
The whole-HUC-8 Peace River domain (\texttt{03100101}; 76 gages) supports per-class analysis.
The NHD-first/SWAT-second method (\S\ref{sec:methods-drainage-audit}; Table~\ref{tab:station-assignment}) mapped 61 gages by exact NHD$\rightarrow$GIS crosswalk, 10 by lake-outlet replacement, and 5 by controlled downstream replacement, yielding 36 clean tributaries, 2 clean mainstem reaches, 16 mainstem reaches with a known NHD offset, and 13 lake-outlet plus 9 canal/artificial cases flagged for review.
A reconciliation funnel reduced these to a calibration-ready set: 54 gages passed the auto-calibration assignment class, 20 had usable 2000--2024 daily streamflow ($\geq$80\% coverage), and 16 were fully calibration-ready.

\begin{table}[htbp]
  \centering
  \caption{NHD-first/SWAT-second USGS station-to-channel assignment for the Peace River HUC-8 domain (Objective~2).}
  \label{tab:station-assignment}
  \small
  \setlength{\tabcolsep}{6pt}
  \renewcommand{\arraystretch}{1.05}
  % Generated by publication/analysis/scripts/emit_tab_station_assignment_tex.py
% Source: /data/SWATGenXApp/codes/web_application/frontend/src/data/stationAssignmentPageCatalog.json
\begin{tabular}{@{}lr@{}}
\toprule
Peace River station assignment (76 USGS gages) & Gages \\
\midrule
\multicolumn{2}{@{}l}{\emph{Assignment class}} \\
\quad Clean tributary & 36 \\
\quad Mainstem, clean & 2 \\
\quad Mainstem, known NHD offset & 16 \\
\quad Lake-outlet review & 13 \\
\quad Canal/artificial review & 9 \\
\midrule
\multicolumn{2}{@{}l}{\emph{Mapping method}} \\
\quad Exact NHD$\rightarrow$GIS crosswalk & 61 \\
\quad Lake-outlet replacement & 10 \\
\quad Controlled downstream replacement & 5 \\
\midrule
\multicolumn{2}{@{}l}{\emph{Calibration reconciliation funnel}} \\
\quad Auto-calibration assignment class & 54 \\
\quad \dots\ with usable 2000--2024 daily flow ($\geq$80\%) & 20 \\
\quad Fully calibration-ready & 16 \\
\bottomrule
\end{tabular}
%
\end{table}

\paragraph{SWAT\textsuperscript{+}--NHD agreement (Peace River).}
All 76 gages matched a SWAT\textsuperscript{+} channel; the median SWAT\textsuperscript{+}/NHD ratio was 1.11 (10th--90th percentile $0.78$--$1.28$), with 66 of 76 inside the loose band and six flagged as assignment outliers (tiny tributary or canal channels picked near a much larger mainstem reach; Figure~\ref{fig:drainage-area-audit-scatter}).

\paragraph{NHD--USGS alignment and joint credibility (Peace River).}
The NHD reference area also tracks the independent USGS NWIS site catalog (Table~\ref{tab:drainage-credibility}): of the 57 gages with a USGS site drainage area, 43 agree within 25\% (15 within 10\%) and the median NHD/NWIS ratio is 0.85---NHD cumulative area on the assigned reach runs about 15\% below the USGS site value, as expected where gages sit on tributaries, lake outlets, or canals whose USGS catalog area is local rather than cumulative; 19 gages lack a USGS site area entirely.
The two alignment questions are largely independent: of the 57 gages with both metrics, 40 (70\%) agree on both, only one (a drainage canal) fails both, and the Pearson correlation between $|$NHD$-$NWIS$|$\% and $|\log(\mathrm{SWAT/NHD})|$ is near zero ($r=-0.005$).
USGS site-file drainage variance therefore does not drive the model--hydrography discrepancies; in particular, all mainstem-offset gages agree with the USGS site area within 25\% while still carrying the SWAT\textsuperscript{+} excess, so that excess is not a catalog-area artifact.

\begin{table}[htbp]
  \centering
  \caption{Peace River two-way drainage-area alignment: NHD reference versus USGS NWIS site area, and joint NHD--USGS / SWAT\textsuperscript{+}--NHD agreement (Objective~2). The two comparisons are near-independent ($r=-0.005$).}
  \label{tab:drainage-credibility}
  \small
  \setlength{\tabcolsep}{6pt}
  \renewcommand{\arraystretch}{1.05}
  % Generated by publication/analysis/scripts/emit_tab_drainage_credibility_tex.py
% Source: /data/SWATGenXApp/codes/web_application/frontend/src/data/stationAssignmentPageCatalog.json
\begin{tabular}{@{}lr@{}}
\toprule
Drainage-area alignment, Peace River & Count / ratio \\
\midrule
\multicolumn{2}{@{}l}{\emph{NHD reference vs USGS NWIS site area} (n=57 with site area; 19 lack it)} \\
\quad NHD within 10\% of NWIS & 15 \\
\quad NHD within 25\% of NWIS & 43 \\
\quad NHD $>$25\% from NWIS & 14 \\
\quad Median NHD/NWIS ratio & 0.85 \\
\midrule
\multicolumn{2}{@{}l}{\emph{Joint NHD--USGS and SWAT--NHD agreement} (n=57)} \\
\quad Both agree ($|$NHD$-$USGS$|\leq$25\%, SWAT/NHD 0.5--2.0$\times$) & 40 \\
\quad NHD--USGS mismatch only & 13 \\
\quad SWAT/NHD mismatch only & 3 \\
\quad Both mismatch & 1 \\
\bottomrule
\end{tabular}
%
\end{table}

\begin{table}[htbp]
  \centering
  \caption{Peace River mainstem drainage-area pipeline trace on a diagnostic panel of representative mainstem reaches: median SWAT\textsuperscript{+}/NHD ratio by stage, localizing the systematic offset (Objective~2). Stage labels defined in Section~\ref{sec:methods-drainage-audit}.}
  \label{tab:drainage-area-pipeline-trace}
  \footnotesize
  \setlength{\tabcolsep}{4pt}
  \renewcommand{\arraystretch}{1.1}
  \resizebox{\textwidth}{!}{%
    % Generated by publication/analysis/scripts/emit_tab_drainage_area_pipeline_trace_tex.py
% Source: /data/SWATGenXApp/codes/publication/analysis/qa/peace-drainage-area-pipeline-trace.csv (mainstem diagnostic-panel gages; n=3)
\begin{tabular}{@{}clcp{5.0cm}@{}}
\toprule
Stage & Quantity & Median ratio vs NHD \texttt{TotDASqKm} & Interpretation \\
\midrule
A & Original NHDPlus HR HU4 \texttt{TotDASqKm} (VAA) & 1.00 & Reference cumulative drainage at the gaged reach \\
B & Original upstream $\Sigma$ \texttt{AreaSqKm} (HU4 network) & 1.00 & Independent cumulative cross-check; agrees with A \\
G & Cleaned flowline reach \texttt{TotDASqKm} & 1.00 & Pre-QSWAT\textsuperscript{+} attribute; $<$0.1\% vs A (VAA not pre-inflated) \\
C & Cleaned flowline-network upstream $\Sigma$ & 0.99 & After topology cleaning; no positive bias \\
D & Post-processed watershed-polygon upstream $\Sigma$ & 1.03 & SWATGenX subbasin polygons; $+2$--$3\%$ only \\
E & QSWAT\textsuperscript{+}/TauDEM channel \texttt{AreaC} at assignment & 1.17 & First large exceedance ($+15$--$17\%$) \\
F & \texttt{chandeg.con} executable area & 1.17 & Headline audit column; matches E \\
H & Live \texttt{rivs1.AreaC} / project SQLite & 1.17 & Matches E/F ($<$0.1\% vs \texttt{chandeg}); offset present before export \\
\bottomrule
\end{tabular}
%
  }
\end{table}

\begin{figure}[htbp]
  \centering
  \includegraphics[width=0.72\textwidth]{final/fig-drainage-area-audit-peace-scatter.png}
  \caption{Peace River HUC-8 drainage-area agreement: SWAT\textsuperscript{+} \texttt{chandeg.con} area versus original NHDPlus HR \texttt{TotDASqKm} at 76 gage-assigned channels (log--log), colored by SWAT\textsuperscript{+}/NHD agreement band; the $1{:}1$ line and the systematic mainstem $+10$--$17\%$ offset are annotated.}
  \label{fig:drainage-area-audit-scatter}
\end{figure}

\paragraph{Where the offset enters.}
The 16 Peace gages in the mainstem-offset class carry a systematic $+10$--$17\%$ excess in SWAT\textsuperscript{+} cumulative area relative to original NHD VAA.
Tracing that area through the pipeline on a diagnostic panel of representative mainstem reaches (Table~\ref{tab:drainage-area-pipeline-trace}) localizes the offset: the original VAA \texttt{TotDASqKm}, the independent upstream catchment-area sums, the pre-QSWAT\textsuperscript{+} \texttt{streams.pkl} reach attribute, and the SWATGenX post-processed watershed polygons all agree with the reference within $\pm3\%$ (median ratios $0.99$--$1.03$).
The excess first appears at the QSWAT\textsuperscript{+}/TauDEM channel-area stage (\texttt{AreaC}, ratio $\approx1.17$) and persists into \texttt{chandeg.con}.
The NHDPlus HR preprocessing therefore preserves cumulative drainage structure; the residual offset is a delineation-definition difference between TauDEM-derived channel contributing area and NHDPlus HR hydrosequence cumulative drainage, introduced downstream of the preprocessing this paper documents.
The live project channel area confirms this directly: \texttt{rivs1.AreaC} and the project SQLite \texttt{gis\_channels} table carry the full mainstem $+17\%$ (Table~\ref{tab:drainage-area-pipeline-trace}, stage~H), and \texttt{chandeg.con} reproduces them to within $0.1\%$---so the SWAT\textsuperscript{+} text export is faithful and the offset originates entirely in the QSWAT\textsuperscript{+}/TauDEM channel-area computation.
These are structural-fidelity diagnostics, not hydrologic validation.

\subsubsection{Objective 3: Structural scalability and model-assembly cost}
\label{sec:results-obj3}

The small, medium, and large benchmark domains (Table~\ref{tab:model-roster}) use the same Model IDs reported in the public runtime benchmark catalog: \texttt{03080102} Oklawaha, Florida; \texttt{09471300} Upper San Pedro, Arizona; and \texttt{03100101} Peace River, Florida (whole HUC8).
No other structural benchmark domains appear in this paper.

\begin{figure}[htbp]
  \centering
  \includegraphics[width=\textwidth]{final/fig-example-basin-maps-combined-3panel.png}
  \caption{Small, medium, and large benchmark domains (size classes S, M, L). Panels (a)--(c) are independently zoomed and are not at a common geographic scale. Structure and model-assembly cost: Table~\ref{tab:benchmark-sm-l-products}.}
  \label{fig:example-basin-maps}
\end{figure}

HRU, channel, and subbasin counts illustrate how geometric and topological workload scales with domain size under fixed national data and discretization settings (473 to 94{,}303 HRUs across the three size classes).
Table~\ref{tab:benchmark-sm-l-products} combines structural discretization, on-disk project size, and measured model-assembly wall time for the same three Model IDs (measured on the production evaluation host; assembly logs archived with the release).
For size classes S and M, assembly time is the end-to-end production worker job duration; for L (whole HUC-8), it is the maximum logged QSWAT\textsuperscript{+} project-assembly phase duration during the documented production build.
Simulation runtime is reported separately in Part~C.

\subsection{Downstream workflow readiness}
\label{sec:results-cal-val-sensitivity}

\subsubsection{Objective 4: Workflow execution on controlled gage basins}
\label{sec:results-obj4}

Two controlled gage basins (\texttt{02297600} and \texttt{05536265}) that are not part of the S/M/L benchmark domain set (Table~\ref{tab:model-roster}) test whether generated projects can enter a standard downstream modeling workflow (Section~\ref{sec:methods-calibration}).
On both basins the platform completed initialization pools, PSO calibration, independent-period forward runs, and Morris batch execution with ensemble diagnostics (Table~\ref{tab:controlled-gage-structure}).
Verification windows were selected separately for each basin based on NWIS data availability and forcing coverage, with no overlap between scored calibration and verification periods (Florida: post-calibration holdout 2019--2024; Illinois: pre-calibration holdout 2012--2015).
These results do not isolate the hydrologic effect of individual preprocessing choices and are not used as proof of national hydrologic validity.

\begin{table}[htbp]
  \centering
  \caption{Structural context for Objective~4 controlled gage basins.}
  \label{tab:controlled-gage-structure}
  \scriptsize
  \setlength{\tabcolsep}{3pt}
  \renewcommand{\arraystretch}{1.05}
  \resizebox{\textwidth}{!}{%
    % Generated by publication/analysis/scripts/emit_tab_controlled_gage_structure_tex.py
\begin{tabular}{@{}llrrrrrrrrllr@{}}
\toprule
Basin ID & State & Ch. & Orig.\ cat. & Orig.\ reach & Div-2 & Isol.\ reach & Final ch. & Final sub. & Cal.\ period & Ver.\ period & HRUs \\
\midrule
\texttt{02297600} & FL & 2 & 41 & 41 & 0 & 4 & 37 & 4 & 2013-2018 & 2019-2024 & --- \\
\texttt{05536265} & IL & 25 & 224 & 224 & 6 & 13 & 206 & 8 & 2020-2024 & 2012-2015 & --- \\
\bottomrule
\end{tabular}
%
  }
\end{table}

\paragraph{Morris sensitivity and ensemble uncertainty.}
Morris analysis (scenario \texttt{Default\_initialized}) used 1000 parameter-space evaluations (27 trajectories, four levels) on each basin's calibration simulation window; streamflow ensembles from 999 completed members provide a 2.5--97.5\% predictive uncertainty band, ensemble median, and best-scoring member trace (Figure~\ref{fig:sensitivity-ensemble}; Table~\ref{tab:sensitivity-ensemble-metrics}).
P-factor and R-factor characterize ensemble machinery on the calibration window, not overall model goodness-of-fit.
Florida's ensemble is narrower (daily P-factor~0.71, R-factor~0.73) than Illinois (daily P-factor~0.87, R-factor~3.13), reflecting a wider 95\% band relative to observed variability on the Illinois calibration window.

\begin{table}[htbp]
  \centering
  \caption{Morris ensemble uncertainty diagnostics (999 completed members; P-factor and R-factor defined in Methods).}
  \label{tab:sensitivity-ensemble-metrics}
  \footnotesize
  \setlength{\tabcolsep}{4pt}
  \renewcommand{\arraystretch}{1.05}
  % Generated by emit_tab_sensitivity_ensemble_metrics_tex.py
% Source: /data/SWATGenXApp/codes/publication/tables/tab-sensitivity-ensemble-metrics.csv
\begin{tabular}{@{}llrrrr@{}}
\toprule
Basin & Step & $N$ & P-factor & R-factor & Best NSE \\
\midrule
\texttt{02297600} & Daily & 999 & 0.705 & 0.733 & 0.817 \\
\texttt{02297600} & Monthly & 999 & 0.722 & 0.894 & 0.882 \\
\texttt{05536265} & Daily & 999 & 0.874 & 3.125 & 0.290 \\
\texttt{05536265} & Monthly & 999 & 0.900 & 5.123 & 0.470 \\
\bottomrule
\end{tabular}
%
\end{table}

\begin{figure}[htbp]
  \centering
  \includegraphics[width=\textwidth,keepaspectratio]{final/fig-sensitivity-ensemble-controlled-basins-daily.png}
  \caption{Morris ensemble daily streamflow for controlled gage basins on each basin's calibration window: (a)~\texttt{02297600} (Florida) and (b)~\texttt{05536265} (Illinois). Shaded band: 2.5--97.5\% of 999 Morris members; dashed gray: ensemble median; navy: observed; dark red: best Morris member. P-factor and R-factor annotated on each panel.}
  \label{fig:sensitivity-ensemble}
\end{figure}

Morris $\mu^*$ indices (Tables~\ref{tab:morris-sensitivity} and~\ref{tab:morris-sensitivity-05536265}) rank parameter influence.
Florida's largest sensitivities include \texttt{surq\_lag}, \texttt{perco}, \texttt{alpha\_bf}, and \texttt{cn3\_swf}; snow parameters show negligible $\mu^*$, consistent with subtropical domain QC.
Illinois shows a contrasting driver set---snow-process parameters (\texttt{melt\_min}, \texttt{melt\_max}) and channel routing (\texttt{k}, \texttt{mann}) among the largest $\mu^*$, together with \texttt{cn3\_swf}, \texttt{perco}, and urban curve-number classes (parameter-ranking spider charts in Supplementary Figure~S\ref{fig:supp-morris-spider}).
Monthly ensembles and calibration/verification hydrographs are in Supplementary Section~\ref{sec:supp-calval-sensitivity}.

\begin{table}[htbp]
  \centering
  \caption{Top eight Morris $\mu^*$ for basin \texttt{02297600} (full 36-parameter ranking in archived \texttt{tab-sensitivity-morris.csv}).}
  \label{tab:morris-sensitivity}
  \footnotesize
  \setlength{\tabcolsep}{4pt}
  \renewcommand{\arraystretch}{1.05}
  % Generated by emit_tab_sensitivity_morris_tex.py
% Source: /data/SWATGenXApp/codes/publication/tables/tab-sensitivity-morris.csv
% site_no: 02297600
\begin{tabular}{@{}rlrrr@{}}
\toprule
Rank & Parameter & $\mu$ & $\mu^*$ & $\sigma$ \\
\midrule
1 & \texttt{surq\_lag} & -0.6435 & 0.6435 & 0.4685 \\
2 & \texttt{perco} & 0.2323 & 0.2880 & 0.4113 \\
3 & \texttt{cn3\_swf} & 0.0598 & 0.1413 & 0.1699 \\
4 & \texttt{alpha\_bf} & -0.0921 & 0.0921 & 0.0987 \\
5 & \texttt{dep\_wt} & 0.0609 & 0.0882 & 0.1565 \\
6 & \texttt{spec\_yld} & 0.0523 & 0.0545 & 0.0945 \\
7 & \texttt{flo\_min} & -0.0142 & 0.0519 & 0.0774 \\
8 & \texttt{dp\_es} & 0.0500 & 0.0500 & 0.0282 \\
\bottomrule
\end{tabular}
%
\end{table}

\begin{table}[htbp]
  \centering
  \caption{Top eight Morris $\mu^*$ for basin \texttt{05536265} (full 36-parameter ranking in archived \texttt{tab-sensitivity-morris.csv}).}
  \label{tab:morris-sensitivity-05536265}
  \footnotesize
  \setlength{\tabcolsep}{4pt}
  \renewcommand{\arraystretch}{1.05}
  \input{../tables/generated/tab-sensitivity-morris-05536265}%
\end{table}

\paragraph{Split-sample calibration diagnostics.}
Full daily and monthly NSE, RMSE (cfs), and PBIAS by workflow stage appear in Supplementary Table~\ref{tab:hydrologic-metrics} (archived \texttt{tab-metrics.csv}).
The summaries below report holdout skill only; they are not interpreted as platform validation.

\subsubsection{Basin 02297600 (Florida)}

NWIS gage \texttt{02297600}; SWAT\textsuperscript{+} routing channel index~2.
Because daily NWIS streamflow for \texttt{02297600} begins on 2009-10-01, the originally planned early calibration window was not feasible.
The final split used 2013--2018 for calibration (three-year warm-up from 2010) and 2019--2024 for independent post-calibration verification (Table~\ref{tab:controlled-gage-structure}).
Scenario \texttt{Default\_calval\_split202606} used a Morris-trimmed parameter subset (36 particles; stagnation after iteration~26).
Run settings are frozen in the archived calibration-settings table (\texttt{tab-calibration-run-settings.csv}; Section~\ref{sec:methods-calibration}).

The initialization pool best captured seasonal timing (daily NSE~0.80).
Calibration improved daily and monthly NSE to 0.84 and 0.90; independent verification retained daily NSE~0.73 and monthly NSE~0.80 with modest wet bias (PBIAS about $+7\%$).
These metrics document workflow execution on a humid subtropical gage basin; they do not establish preprocessing optimality or national skill.

\subsubsection{Basin 05536265 (Illinois)}
\label{sec:results-cal-val-05536265}

NWIS gage \texttt{05536265}; SWAT\textsuperscript{+} routing channel index~25.
Calibration used scored window 2020--2024 after a two-year warm-up from 2018.
Verification used an independent pre-calibration holdout on 2012--2015 after a one-year warm-up from 2011.
Scenario \texttt{Default\_calval\_split202606} used a Morris-trimmed subset (48 particles; stagnation after iteration~26).
The initialization pool best was weak on the calibration window (daily NSE~0.20).
Calibration global best reached daily NSE~0.37 and monthly NSE~0.60; verification skill remained modest (daily NSE~0.22, monthly NSE~0.42).
Illinois reports low holdout skill as a workflow diagnostic on a second hydroclimate; automated assembly does not ensure calibration-ready hydrology, and process causes for weak skill are not diagnosed here.
Morris screening still completed (999 members) with the wider ensemble band noted above, illustrating that workflow execution does not replace expert process review and calibration.

\subsection{Operational simulation cost}
\label{sec:results-runtime-benchmark}

\subsubsection{Objective 5: SWAT\textsuperscript{+} simulation runtime benchmark}
\label{sec:results-obj5}

Objective~5 reports measured SWAT\textsuperscript{+} simulation cost on one documented host/build configuration as supporting evidence for computational feasibility (Table~\ref{tab:model-roster}).
Primary results use the small, medium, and large size classes; the six-point scaling ladder (X20/X40/X60) is summarized in Supplementary Section~\ref{sec:supp-runtime-benchmark}.
Extended scenarios and the full six-point ladder are on the public benchmark page (\url{https://www.swatgenx.com/swat-plus-runtime-benchmark}).

\paragraph{Measurement protocol.}
All runs use calendar 2021 (365 days) on AMD EPYC 7282 (10~vCPU).
Production binary: Intel ifx \texttt{-O3 -ipo}.
Primary profile: filtered NetCDF \texttt{channel\_sd} at gauge channels via \texttt{print\_filter.prt}.

\paragraph{Primary calibration-profile scaling (S/M/L).}
Table~\ref{tab:runtime-benchmark}: filtered runs scale from about 5~s (473 HRUs) to about 100~s (11{,}284 HRUs) and about 23~min (94{,}303 HRUs) per simulated year.
Print-scope, export-format, compiler, and six-point ladder tables appear in Supplementary Section~\ref{sec:supp-runtime-benchmark}.

\begin{table}[htbp]
  \centering
  \caption{One-year simulation wall time under calibration filtered NetCDF profile for the small, medium, and large benchmark domains.}
  \label{tab:runtime-benchmark}
  \small
  % Generated by emit_tab_runtime_benchmark_tex.py
% Source: /data/SWATGenXApp/codes/publication/tables/tab-runtime-benchmark.csv
\begin{tabular*}{\textwidth}{@{\extracolsep{\fill}}llrrrr@{}}
\toprule
Size class & Model & HRUs & Wall (s) & s/day & Peak RSS (KB) \\
\midrule
S & \texttt{03080102} & 473 & 5.42 & 0.01 & 55924 \\
M & \texttt{09471300} & 11284 & 99.79 & 0.26 & 511468 \\
L & \texttt{03100101} & 94303 & 1378.13 & 3.37 & 3956648 \\
\bottomrule
\end{tabular*}
%
\end{table}

\begin{figure}[htbp]
  \centering
  \includegraphics[width=0.78\textwidth]{final/fig-runtime-benchmark-hru-scaling.png}
  \caption{Wall time versus HRU count for the six published benchmark domains (Supplementary Table~\ref{tab:runtime-benchmark-hru-scaling}).}
  \label{fig:runtime-benchmark-hru-scaling}
\end{figure}
